\documentclass[10pt,letterpaper]{article}
\usepackage[margin=0.65in]{geometry}
\usepackage[T1]{fontenc}\usepackage[utf8]{inputenc}\usepackage{charter}
\usepackage{xcolor}\usepackage[normalem]{ulem}\usepackage{hyperref}
\newcommand{\del}[1]{{\color{red}\sout{#1}}}
\newcommand{\add}[1]{{\color{blue}\uline{#1}}}
\setlength{\parindent}{0pt}\setlength{\parskip}{6pt}
\emergencystretch=3em
\begin{document}
\textbf{Ziyu Xu -- Full-text redline against original resume}\\
\del{Red strikeout: deleted} \quad \add{Blue underline: added}\\
Ziyu\_Xu\_master\par
Original baseline: CMU\_Ziyu\_original.tex. All visible text is compared in PDF reading order; moved content appears as deletion and insertion. Layout-only changes are described in changes.md.
\hrule\medskip
Ziyu (Rick) Xu Robotics Institute, Carnegie Mellon University Research Interests: Robot Foundation Models, Real-World Policy Adaptation, and Learning-Augmented Control +1 7348005258 | rickx@andrew.cmu.edu \par\medskip\textbf{Education}\par  Carnegie Mellon University (CMU) Aug. 2026 \add{--} \add{Expected} \add{2028} Master of Science in Robotics (MSR) Pittsburgh, PA University of Michigan, Ann Arbor (UM) Oct. 2024 -- May 2026 B.S.E. in Mechanical Engineering (Robotics Track); Minor in Computer Science GPA: 3.85/4.0 (Top 5\%)  Core Coursework: Robot Learning, RL, Embedded Systems, Optimal Control, Convex Optimization Shanghai Jiao Tong University (SJTU) Sep. 2022 -- Jul. 2026 B.S.E. in Electrical and Computer Engineering GPA: 3.56/4.0 (Top 20\%)  Core Coursework: Algorithms, Mathematical Foundations, Probability Theory, Signals and Systems \par\medskip\textbf{Research Experiences and Internships}\par  \add{Intelligent} \add{Control} \add{Lab,} \add{Carnegie} \add{Mellon} \add{University} \add{Aug.} \add{2026} \add{--} \add{Present} \add{Graduate} \add{Student} \add{Researcher,} \add{advised} \add{by} \add{Prof.} \add{Changliu} \add{Liu} \add{Pittsburgh,} \add{PA} \add{} \add{Whole-Body} \add{and} \add{Safe} \add{Locomotion} \add{Research} \add{(ongoing):} \add{Investigating} \add{whole-body} \add{control,} \add{safety} \add{with} \add{control} \add{barrier} \add{functions} \add{(CBFs),} \add{and} \add{world-model-based} \add{locomotion.} \add{} \add{Dexmate} \add{Bimanual} \add{Manipulation} \add{Demo,} \add{IROS} \add{2026:} \add{Contributed} \add{to} \add{a} \add{bimanual} \add{block-stacking} \add{demonstra-} \add{tion} \add{combining} \add{agent} \add{planning} \add{with} \add{policy} \add{execution.} \add{Implemented} \add{the} \add{agent} \add{component,} \add{prepared} \add{the} \add{on-site} \add{environment,} \add{adapted} \add{the} \add{system} \add{to} \add{different} \add{grippers,} \add{and} \add{performed} \add{system} \add{calibration.} Unitree Robotics, R\&D Department, Embodied AI Team May 2026 -- Present Research Intern, Humanoid Robot Foundation Models Hangzhou, China  Closed-Loop Real-Robot Infrastructure for Humanoid Foundation-Model Policies: Built and validated an end- to-end policy infrastructure on Unitree H2/G1 humanoids, integrating teleoperation data collection, VLA/WAM and whole-body policy adapters, rollout logging and evaluation, controller interfaces, and policy fine-tuning. Enabled launch with one command and closed-loop rollouts across platforms, allowing policies and controllers to be integrated without modifying platform-specific launch workflows.  Embodiment VLA Adaptation \& Post-Training: Conducted a structured survey of representative robot VLA, WAM, and physical-AI policies, especially work from NVIDIA GEAR Lab. Deployed and fine-tuned Isaac GR00T/OpenPI VLA policies with SONIC whole-body control on Unitree humanoids using real-robot rollout traces. Currently developing a phase-aware adaptation method that aligns semantic task progress and execution timing across embodiments for policy transfer and rollout-driven post-training.  Reproducible Real-World VLA Benchmark for Humanoid Manipulation: Implemented an 8-task, 24-scene benchmark with green-screen background replacement to evaluate foundation models under varied visual con- ditions. Built a PICO-based scene-digitization and AR-guided reconstruction pipeline that captured shelf-object layouts as reusable digital scene templates, and registered target placements back into the physical workspace, achieving 2.5 cm median placement error. Shanghai Artificial Intelligence Lab \& IRMV, SJTU Apr. 2026 -- Present Research Assistant, Directed by Prof. Junchi Yan and Prof. Hesheng Wang Shanghai, China  Context-Aware Re-encoding for Iterative Visual KV Compression: Identified representation errors caused by reusing retained visual KV entries after their spatial neighbors are removed, and showed that these errors com- pound across repeated compression in multi-turn and streaming inference. Developed a re-encoding method that recomputes each retained set in its new context at the original positions, and benchmarked it against inheritance and restoration methods under matched memory and runtime budgets. Computational Autonomy and Robotics Lab, UM Jun. 2025 -- Apr. 2026 Research Assistant, Directed by Prof. Maani Ghaffari Ann Arbor, MI  Online Learning Lie-MPC for Robust ASV Control: Developed a Neural-Fly-inspired learning-augmented Lie-MPC framework for an autonomous surface vehicle, estimating residual wind/current disturbances online and injecting learned dynamics corrections into the controller for robust trajectory tracking under hydrodynamic uncertainty. Implemented and evaluated the framework on a real ASV, achieved a 20\% improvement in trajectory-tracking performance over the baseline MPC controller during field experiments.  Field Perception for Aquatic Vegetation Tracking: Built a water-scene perception pipeline for autonomous sur- face vehicles, combining DINO/SAM2-based segmentation with Grounded-DINO and YOLO fine-tuning to detect submerged vegetation and provide perception signals for navigation and autonomy experiments.  ASV Hardware and Sensor Integration for Closed-Loop Field Autonomy: Designed sensor mounts and inte- grated GPS, lidar, IMU, and camera systems on an autonomous surface vehicle. Aligned and calibrated sensors to support field data collection and control experiments. Intelligent Robotics and Autonomy Lab , UM Oct. 2024 -- Apr. 2026 Research Assistant, Directed by Prof. Vasileios Tzoumas Ann Arbor, MI  Multi-Robot Active-SLAM Simulation \& ROS2 Integration: Rebuilt a Unity + AirSim multi-drone simulation testbed with redesigned vehicle dynamics and sensing modules, then migrated SlideSLAM communication logics from ROS1 to ROS2 by refactoring topic interfaces and workflows for bandwidth-limited active-SLAM evaluation.  Resource-Aware Multi-Robot Exploration Algorithms: Developed distributed exploration EXP3-style algorithms for bandwidth-limited robot teams, formulating communication-aware action selection with submodular objec- tives and bandit feedback in simulation to balance map coverage and communication cost. Surgical Computer Vision Detection, SJTU Apr. 2024 -- Dec. 2024 Research Assistant, Directed by Prof. Yutong Ban Shanghai, China  Foundation Segmentation Model Evaluation for Surgical Video: Evaluated Segment Anything Model variants on surgical video frames, analyzing domain shift and temporal mask stability around instruments, tissues, and tool-tissue interaction regions. \par\medskip\textbf{Publications and Manuscripts}\par   Xiaohan Wang* , Ziyu Xu* , Xuyi Yang. ``Re-encoding What You Keep After Every Compression: Retained Visual KV Is Computed in a Context That No Longer Exists'', Submitted to ICLR 2027. * Equal contribution.  Xuyi Yang, Xiaohan Wang, Wenhao Zhang, and Ziyu Xu. ``StreamJIT: Continuous Observation and Just-in-Time Deliberation for Proactive Streaming Video Assis- tants'', Submitted to ICLR 2027.  Yinan Dong, Ziyu Xu, Tsimafei Lazouski, Sangli Teng, and Maani Ghaffari. ``Online Learning-Enhanced Lie Algebraic MPC for Robust Autonomous Surface Vehicle Control'', IEEE Robotics and Automation Letters (RA-L), revised manuscript under review, 2025.  C. Yuan, J. Jiang, K. Yang, Z. Xu, . . . , and Y. Ban. ``Systematic Evaluation and Guidelines for Segment Anything Model in Surgical Video Analysis'', npj Digital Surgery, accepted, 2024. \par\medskip\textbf{Awards}\par   Gold Medal in 2023 University Physics Competition, by American Physical Society \& American Astronomi- cal Society (top 1\%)  Third place in 2023 RoboMaster Campus Competition (top 10\%)  Dean's Honor List, University of Michigan College of Engineering (2024, 2025)  M Prize in 2024 Mathematical Contest in Modeling (top 13\%) \par\medskip\textbf{Skills}\par   Robot Learning \& Control: VLA post-training, humanoid foundation models, whole-body policy, RL, real- robot policy deployment, teleoperation/data collection, MPC, sim-to-real  Perception \& Multimodal ML: PyTorch, OpenCV, SAM2, Grounded-DINO, YOLO, DINO features, visual token compression, spatial grounding, LLMs fine-tuning  Systems \& Tools: ROS/ROS2, Python/C++, Unity + AirSim, Docker, Linux, Git, sensor integration, CAD, FEA, Manufacturing, embedded system\end{document}